\cotefile{156-4}
\scanpage{74}\mspage{73}

\deleted{\uline{Corollaire} \unsure{2} \unsure{La}}

\uline{\unsure{Définition|Definition}} On dit \deleted{\unsure{alors}} \inserted{que $G$ \uline{raffine} $F$} \deleted{\unsure{si|dit} \unsure{est un} \unsure{raffint|raffinement}} ou est un \uline{raffinement} de $F$ \deleted{\unsure{.} On écrit \unsure{alors}} si $\deleted{\unsure{\in}}\exists$ un $F'$ avec \deleted{$G\ll F$} $G\leq F'\preceq F$. On écrit $G\ll F$

\uline{Corollaire} \deleted{\unsure{2}} \unsure{la|La} relation $G\ll F$ est une relation d'ordre (impliquée par $G\leq F$ et par $G\preceq F$)
\margin{\uline{\unsure{Utilise|Utiliser} $C6$,} \unsure{ni} \uline{$C5$ ni $C7$}.}

Transitivité : Soit $F\deleted{\unsure{\leq}}\ll G\ll H$, donc

\diagram{\begin{tikzcd}[ampersand replacement=\&,row sep=0.1em,column sep=0.3em]
 {} \& \deleted{G} \& {} \& \deleted{H'} \& \deleted{\preceq} \& \deleted{H} \\
 {} \& {} \& {} \& \deleted{\rotatebox{90}{$\leq$}} \& {} \& {} \\
 \deleted{H} \& \deleted{G'} \& \deleted{\preceq} \& \deleted{G} \& {} \& {} \\
 {} \& \deleted{\rotatebox{90}{$\leq$}} \& {} \& {} \& {} \& {} \\
 {} \& \deleted{\unsure{H|G}}\deleted{F} \& {} \& {} \& {} \& {}
 \arrow[from=1-2,to=3-1,dash]\arrow[from=1-4,to=4-2,dash]
 \arrow[from=1-5,to=3-3,dash]\arrow[from=1-6,to=3-5,dash]
\end{tikzcd}}

\diagram{\begin{tikzcd}[ampersand replacement=\&,row sep=0.08em,column sep=0.22em]
 H \& \succeq \& H' \& {} \& {} \\
 \vdots \& {} \& \rotatebox{90}{$\leq$} \& {} \& {} \\
 {} \& \cdots \& G \& \succeq \& G' \\
 {} \& {} \& {} \& {} \& \rotatebox{90}{$\leq$} \\
 {} \& {} \& {} \& {} \& \deleted{\unsure{H}}F
\end{tikzcd}}

on complète par $\unsure{6\text{b}|6\text{a}|C6|C5}$

\diagram{\begin{tikzcd}[ampersand replacement=\&,row sep=0.08em,column sep=0.22em]
 H \& \succeq \& H' \& {} \& {} \& {} \\
 \rotatebox{90}{$\leq$} \& \circ \& \rotatebox{90}{$\leq$} \& {} \& {} \& {} \\
 \bar H \& \succeq \& G \& \succeq \& G' \& \quad (*) \\
 \rotatebox{90}{$\leq$} \& \circ \& \rotatebox{90}{$\leq$} \& \circ \& \rotatebox{90}{$\leq$} \& {} \\
 K \& \succeq \& \bar G \& \succeq \& F \& {}
\end{tikzcd}}

et on gagne, \unsure{car|on|ou}
\[
 F\preceq K\leq H
\]

Supposons maintenant que $F=H$ i.e. $F\ll G$ et $G\ll F$, prouvons $F=\deleted{\unsure{H}}G$. On \deleted{\unsure{tire de}}
\[
 \begin{aligned}
 F\preceq K\leq{} & F & \text{on va en déduire}\\[-2pt]
 &\rotatebox{90}{$= $}\\[-2pt]
 &H & K=F
 \end{aligned}
\]
et comme $\underset{\substack{\rotatebox{90}{$=$}\\F}}{K}\leq\bar H\leq F$ \unsure{etc} on obtiendra
\[
 \bar H=F\quad \deleted{\text{et de m\textsuperscript{me}}\
 \unsure{\bar G|\bar H}=\unsure{F|K}\ (\text{comme } F\preceq\bar G\preceq
 \underset{\substack{\rotatebox{90}{$=$}\\F}}{K})\ \bar G=F}
\]
donc $G\preceq F$. De façon symétrique on aura $F\preceq G$, d'où $F=G$. Donc il reste à prouver le

\uline{Lemme} Soient $F,K\in\unsure{\mathcal F|\mathscr F}$ avec $F\preceq K\leq F$. Alors $K=F$ (i.e. si $F$ est une subdivision d'une sous-figure $K$, alors $F=K$)
